% !TEX root = TFM-ML-01-memoria-referencia.tex
\section{Criterios de diseño del servicio}\label{sec:estrategias}\label{cap:criterios}


\subsection{Requisitos del servicio (QFD)}\label{sec:qfd}

El Despliegue de la Función de Calidad (\textit{Quality Function Deployment},
QFD) traduce las necesidades del usuario en requisitos técnicos del servicio.
Su estructura matricial permite priorizar los parámetros de diseño en función
de su capacidad para satisfacer las necesidades de mayor peso relativo y,
simultáneamente, evaluar la posición competitiva del servicio propuesto
frente a las alternativas existentes en el corredor.

\subsubsection{Estructura de la matriz}
\label{sec:qfd-estructura}

La matriz elaborada incorpora dieciséis requerimientos del cliente (filas)
extraídos de los cinco insights de diseño (§\ref{sec:insights-diseno}) y
dieciséis requerimientos funcionales (columnas) que constituyen los
parámetros técnicos del servicio. Cada requerimiento funcional incluye
un objetivo de mejora medible.

Los requerimientos del cliente se ponderan en función de la frecuencia e
intensidad de los problemas registrados en la encuesta. Los cinco con
mayor peso relativo (8,3\,\% cada uno) son: la regularidad confiable,
el uso sin reserva previa, la cobertura desde el municipio de origen,
la viabilidad económica del desplazamiento y el tiempo de descanso real
tras el turno \parencite{LouridoRegueira2026}.

La evaluación competitiva compara el servicio propuesto frente a cuatro
operadores de referencia en el corredor: Tranvías de Ferrol, Monbus,
Renfe Cercanías C-1 y Cabify. Ninguno opera en la franja 22:00--08:00
de forma específica para personal sanitario, lo que sitúa al servicio
propuesto en una posición diferencial sin competidor directo
\parencite{Alsaferrol2024horarios,RenfeCercaniasFerrol2026}.

\subsubsection{Resultados de la matriz}
\label{sec:qfd-resultados}

El análisis de la importancia técnica relativa identifica cuatro
requerimientos funcionales con mayor capacidad de satisfacción agregada,
todos ellos con puntuación máxima de 5,0 o próxima a ella.

En primer lugar, la \textbf{disponibilidad de espacio de espera cubierto
e iluminado} en los nodos P+R (objetivo: 100\,\% de paradas con cubierta
e iluminación nocturna) obtiene la mayor puntuación técnica, con incidencia
directa en los requerimientos de seguridad percibida y espacio de espera
digno. Este resultado refuerza el criterio derivado del Insight~5: la
percepción de seguridad en franja nocturna no puede delegarse en el
vehículo si el espacio de espera previo es inseguro.

En segundo lugar, la \textbf{información en tiempo real del vehículo}
(objetivo: actualización de posición cada 30\,segundos o menos) resuelve
de forma transversal los requerimientos de previsibilidad y reducción de
la incertidumbre en la espera nocturna. Trazable a Insight~3.

En tercer lugar, la \textbf{capacidad garantizada sin sobreocupación}
(objetivo: 0\,\% de usuarios sin plaza en la franja crítica post-turno)
es el parámetro determinante para la regularidad percibida. Un servicio
con plazas no garantizadas no puede resolver el patrón documentado en
el Insight~1. Trazable a Insight~1.

Finalmente, el \textbf{interior confortable para el trayecto post-turno}
(objetivo: satisfacción igual o superior a 4 sobre 5 en encuesta de
usuario) responde directamente a la necesidad de descanso activo durante
el desplazamiento, identificada como componente del riesgo de fatiga
\parencite{LouridoRegueira2026}.

\begin{figure}[htbp]
	\centering
	\fbox{\rule{0pt}{6cm}\rule{0.95\textwidth}{0pt}}
	% Sustituir por: \includegraphics[width=\textwidth]{../03-graficas/TFM-ML-GFX-qfd.pdf}
	\caption{Matriz QFD del servicio de movilidad nocturna CHUF--Ferrol.
		Dieciséis requerimientos del cliente ponderados frente a dieciséis
		requerimientos funcionales con objetivos medibles. Evaluación competitiva
		frente a Tranvías de Ferrol, Monbus, Renfe Cercanías C-1 y Cabify.
		Elaboración propia \parencite{LouridoRegueira2026}.}
	\label{fig:qfd}
\end{figure}

Los resultados de la matriz orientan directamente la priorización
MoSCoW (§\ref{sec:moscow}): los requerimientos funcionales con mayor
importancia técnica pasan a la categoría \textit{Must~have}, mientras
que los de importancia media y baja se distribuyen entre
\textit{Should~have} y \textit{Could~have} en función de su
viabilidad operativa en el horizonte 2035.

\subsection{Priorización MoSCoW}\label{sec:moscow}

La priorización MoSCoW organiza los requerimientos funcionales
identificados en la matriz QFD (§\ref{sec:qfd}) en cuatro categorías
según su carácter imprescindible para el funcionamiento del servicio.
El criterio de clasificación combina el peso relativo obtenido en la
matriz con la trazabilidad directa a los insights de diseño: un
requerimiento funcional es \textit{Must~have} si su ausencia impide
que el servicio resuelva alguno de los cinco problemas documentados
empíricamente \parencite{LouridoRegueira2026}.

\subsubsection{\textit{Must have} - Imprescindible}
\label{sec:moscow-must}

Los requerimientos siguientes son condición necesaria del servicio.
Su ausencia lo invalida como solución al problema documentado.

\begin{itemize}

	\item \textbf{Horario sincronizado con los cambios de turno del CHUF
	(22:00 y 08:00).} El 52,1\,\% de la muestra carece de transporte
	público fiable en la franja nocturna \parencite{LouridoRegueira2026}.
	Un servicio con frecuencia genérica no resuelve esta necesidad.
	Objetivo: 100\,\% de los cambios de turno cubiertos.

	\item \textbf{Uso sin reserva previa, por horario fijo.}
	La demanda espontánea de shuttle formulada por los usuarios
	\parencite{LouridoRegueira2026} implica acceso inmediato, no
	planificación anticipada. Un sistema de reserva obligatoria
	reproduce la barrera que el servicio pretende eliminar.
	Objetivo: 0\,\% de trayectos con reserva obligatoria.

	\item \textbf{Capacidad garantizada en la franja crítica post-turno.}
	El 52,7\,\% de los encuestados conduce con fatiga varias veces a la
	semana \parencite{LouridoRegueira2026}. Un servicio con plaza no
	garantizada obliga al trabajador a mantener el vehículo privado como
	alternativa, anulando el efecto del servicio.
	Objetivo: 0\,\% de usuarios sin plaza en la franja 07:30--09:00.

	\item \textbf{Espacio de espera cubierto e iluminado en todos los
	nodos P+R.} Requerimiento con mayor importancia técnica en la
	matriz QFD. La percepción de seguridad de las usuarias que viajan
	solas en franja nocturna \parencite{LouridoRegueira2026} no puede
	garantizarse únicamente desde el interior del vehículo si el espacio
	de espera previo es inseguro.
	Objetivo: 100\,\% de paradas con cubierta e iluminación nocturna.

	\item \textbf{Recorrido directo sin transbordos.}
	El servicio se articula como transporte punto a punto entre nodos
	P+R y el hospital \parencite{LouridoRegueira2026}. Un transbordo
	intermedio añade tiempo, incertidumbre y exposición en franja
	nocturna, contradiciendo los criterios de los Insights 1 y 5.
	Objetivo: 0 transbordos en el 100\,\% de los trayectos.

	\item \textbf{Cobertura interurbana hacia los municipios del área.}
	El 38,9\,\% de los encuestados reside en un municipio distinto al
	del hospital \parencite{LouridoRegueira2026}. Los nodos P+R deben
	situarse en los accesos naturales al corredor, no solo en el núcleo
	urbano de Ferrol. Objetivo: cobertura de al menos el 80\,\% de los
	municipios de origen de los trabajadores del CHUF.

	\item \textbf{Tiempo de espera máximo de 15~min desde el fichaje
	de salida.} La fatiga post-turno condiciona el umbral de tolerancia
	de espera: superado ese umbral, el trabajador descarta el servicio
	y recurre al vehículo privado o al taxi (Participante~1, entrevista
	semiestructurada, 22~de mayo de 2026). Trazable a Insight~1 e Insight~3.
	Objetivo: tiempo de espera $\leq$\,15~min en el 100\,\% de los
	servicios post-turno. Fuente cualitativa (n=1); requiere validación
	con muestra representativa para fijarse como umbral técnico.

\end{itemize}

\subsubsection{\textit{Should have} - Necesario para la adopción}
\label{sec:moscow-should}

Los requerimientos siguientes aumentan significativamente la adopción
y la calidad del servicio. Su ausencia no invalida la solución, pero
reduce su eficacia.

\begin{itemize}

	\item \textbf{Interior confortable para el trayecto post-turno.}
	El trayecto de retorno tras 10--12\,h de turno nocturno es tiempo
	de descanso activo \parencite{LouridoRegueira2026}. El diseño
	interior debe facilitarlo, no impedirlo.
	Objetivo: satisfacción $\geq$\,4/5 en encuesta de usuario.

	\item \textbf{Información en tiempo real del vehículo.}
	Reduce la incertidumbre en la espera nocturna y refuerza la
	percepción de control del usuario sobre su desplazamiento.
	Trazable a Insight~3.
	Objetivo: actualización de posición cada $\leq$\,30\,segundos.

	\item \textbf{Tarifa plana para personal sanitario.}
	El coste del desplazamiento nocturno es una barrera documentada
	\parencite{LouridoRegueira2026}. Una tarifa diferenciada para
	trabajadores del CHUF actúa como incentivo de adopción y como
	reconocimiento institucional del servicio.
	Objetivo: reducción del coste de desplazamiento nocturno en al
	menos un 70\,\% respecto al taxi.

	\item \textbf{Frecuencia garantizada en festivos y fines de semana.}
	El 81,4\,\% de la muestra trabaja en régimen rotatorio con noches
	\parencite{LouridoRegueira2026}, lo que incluye festivos. Un
	servicio que cesa en días no laborables excluye a la mayoría de
	los usuarios en sus noches de mayor necesidad.
	Objetivo: 100\,\% de servicios operativos en festivos.

	\item \textbf{Gobernanza del servicio con participación del hospital.}
	La credibilidad operacional ante los usuarios depende de que el CHUF
	participe en la gestión: el hospital conoce los turnos, la variabilidad
	de salida por planta y las excepciones operativas (Participante~1,
	entrevista semiestructurada, 22~de mayo de 2026). Trazable a Insight~4.
	El CHUF o SERGAS debe actuar como socio clave del servicio, no solo
	como destinatario o regulador.
	Objetivo: integración progresiva de representantes del \textsc{chuf}
	en el comité de seguimiento del servicio durante el primer año
	de operación.

\end{itemize}

\subsubsection{\textit{Could have} - Valor diferencial}
\label{sec:moscow-could}

Los requerimientos siguientes aportan valor añadido y refuerzan
la identidad del servicio, pero no son críticos para su viabilidad
en la fase inicial de implantación.

\begin{itemize}

	\item \textbf{Identidad visual reconocible entre los trabajadores
	del CHUF.} Contribuye a la adopción y al sentimiento de
	pertenencia al servicio. No resuelve el problema de transporte
	pero refuerza su sostenibilidad social.
	Objetivo: reconocimiento de marca $\geq$\,80\,\% entre la plantilla.

	\item \textbf{Reducción de emisiones respecto al vehículo privado.}
	El vehículo eléctrico autónomo como candidato tecnológico preferente
	implica esta ventaja de forma inherente \parencite{SHOW2024Legacy}.
	Objetivo: reduction de emisiones $\geq$\,50\,\% respecto al coche
	particular.

	\item \textbf{Canal de aviso de incidencias en tiempo real.}
	Mejora la experiencia ante interrupciones del servicio y reduce
	la ansiedad asociada a la dependencia de un único medio de
	transporte nocturno. Trazable a Insight~3.
	Objetivo: notificación en menos de 2\,minutos desde la detección.

\end{itemize}

\subsubsection{\textit{Won't have} - Fuera del alcance}
\label{sec:moscow-wont}

Los requerimientos siguientes quedan explícitamente fuera del
alcance del servicio en el horizonte de diseño 2035.

\begin{itemize}

	\item \textbf{Transporte puerta a puerta al domicilio de cada
	trabajador.} Operativamente inviable: la dispersión residencial
	de la plantilla hace imposible un servicio de recogida individual
	que sea puntual y económicamente sostenible. El servicio conecta
	nodos P+R con el hospital, no domicilios con el hospital.

	\item \textbf{Cobertura del 100\,\% de los municipios del área
	sanitaria desde el primer día.} La implantación progresiva por
	corredores de mayor demanda es la estrategia operativa coherente
	con el modelo P+R. La cobertura total es un objetivo de escala,
	no de diseño inicial.

\end{itemize}

\begin{figure}[htbp]
	\centering
	\fbox{\rule{0pt}{5cm}\rule{0.95\textwidth}{0pt}}
	% Sustituir por: \includegraphics[width=\textwidth]{../03-graficas/TFM-ML-GFX-moscow.pdf}
	\caption{Priorización MoSCoW de los requerimientos funcionales
		del servicio de movilidad nocturna CHUF--Ferrol. Elaboración
		propia a partir de los resultados de la matriz QFD.}
	\label{fig:moscow}
\end{figure}

\subsection{Valores del servicio}\label{sec:valores}

Los valores del servicio constituyen la fundamentación conceptual
que orienta cada decisión proyectual posterior. No son atributos
técnicos (recogidos en el QFD, §\ref{sec:qfd}), sino el marco
de significado desde el que se diseña: el porqué de cada criterio,
materializado en tres categorías que operan de forma simultánea
sobre la propuesta.

\subsubsection{Valores racionales}
\label{sec:valores-racionales}

El servicio responde a una necesidad documentada empíricamente: el
52,7\,\% del personal sanitario nocturno del área de A Coruña
conduce con fatiga varias veces a la semana porque no existe
alternativa operativa en la franja 22:00--08:00
\parencite{LouridoRegueira2026}. El valor racional central es la
\textbf{reducción del riesgo}: eliminar el desplazamiento en
vehículo privado bajo fatiga post-turno. A este se suman la
\textbf{eficiencia temporal} (horarios sincronizados con los
cambios de turno, sin tiempo de espera imprevisible) y la
\textbf{viabilidad económica} (una tarifa plana para personal
sanitario significativamente inferior al coste del taxi o del
aparcamiento privado hospitalario).

\subsubsection{Valores experienciales}
\label{sec:valores-experienciales}

El trayecto de vuelta a casa tras diez o doce horas de turno
nocturno no es un desplazamiento neutro: es el primer momento
de descanso disponible. El servicio debe \textbf{proteger ese
tiempo}, no consumirlo en incertidumbre, incomodidad o alerta.
El valor experiencial prioritario es la \textbf{seguridad
percibida}: la usuaria que viaja sola a las 08:00 debe sentirse
tan segura en el nodo de espera como dentro del vehículo. La
brecha de 37 puntos porcentuales en riesgo de fatiga entre
mujeres (90,8\,\%) y hombres (53,8\,\%) \parencite{LouridoRegueira2026}
convierte este valor en un criterio de diseño no negociable.
A él se añaden la \textbf{previsibilidad} (saber exactamente
cuándo llega el vehículo) y el \textbf{confort activo} (un
interior que permita cerrar los ojos o descansar durante los
veinte minutos del trayecto).

\subsubsection{Valores culturales y simbólicos}
\label{sec:valores-simbolicos}

El sistema sanitario público funciona porque hay personas que
trabajan de noche. El servicio propuesto tiene un valor simbólico
que trasciende el transporte: es el \textbf{reconocimiento
institucional} de que ese trabajo nocturno merece una
infraestructura específica. Este valor conecta el servicio con
una lectura más amplia de \textbf{equidad de género en la
movilidad}: en una muestra 91\,\% femenina, diseñar transporte
nocturno seguro es diseñar para las mujeres
\parencite{LouridoRegueira2026}. Por último, la electrificación
del corredor y la reducción de vehículos privados en el acceso
al CHUF aportan un valor de \textbf{sostenibilidad colectiva}
coherente con el horizonte normativo europeo
\parencite{CE_Movilidad2020}.

\subsection{Grado de innovación del servicio propuesto}
\label{sec:grado-innovacion}
\begin{pendiente}
\end{pendiente}

